\chapter{Designing a learning API}
\label{ch:learning-api}

\begin{learningobjectives}
After this chapter, you should be able to:
\begin{itemize}
  \item identify the mathematical structure shared by linear regression,
  logistic regression, and a linear support vector machine;
  \item separate an estimator interface from a differentiable component
  interface;
  \item explain why activation and loss choices form a coupled contract;
  \item specify shape, state, numerical, and randomness guarantees; and
  \item test an analytical gradient independently with finite differences.
\end{itemize}
\end{learningobjectives}

\begin{chapterconnection}
Chapters~\ref{ch:linear-regression}--\ref{ch:linear-svm} developed three
single-neuron configurations: linear regression, logistic regression, and a
linear support vector machine. Their formulas differ in meaning, but their
implementations repeat an affine score, parameter updates, validation, and
fitted-state handling. This chapter asks which repetitions express genuine
structure and which encode different scientific claims. The resulting
components prepare us to build a multilayer network without treating hidden
neurons as independently fitted estimators.
\end{chapterconnection}

\section{An API is a scientific boundary}

An application programming interface, or API, is the documented boundary
through which one piece of software uses another. For a learned model, the
boundary includes more than method names. It states what one row represents,
which shapes are accepted, what fitting changes, what an output means, which
invalid states are rejected, and which numerical behavior callers may rely on.

Suppose $X\in\mathbb R^{n\times p}$ contains $n$ observations and $p$
features. A single affine unit calculates
\[
  z=Xw+b\mathbf 1,
  \qquad
  a=\phi(z).
\]
Here $z$ is the \emph{preactivation}, $a$ the \emph{postactivation}, $w$ a
length-$p$ weight vector, and $b$ a scalar bias. This computation alone does
not determine a statistical model. The activation, target domain, loss,
sampling assumptions, and prediction policy supply meaning.

\begin{definitionbox}{Interface contract}
An interface contract records the conditions a caller must satisfy, the state
an operation may change, the result and semantics it guarantees, and the
failures it reports when it cannot honor that guarantee.
\end{definitionbox}

For example, a binary probability estimator may require labels in
$\{0,1\}$ and return values in $(0,1)$. A label decision requires an additional
threshold and cost policy. Returning a number in $(0,1)$ is not by itself
evidence of calibration.

\conceptcheck{Two classes both expose methods named \texttt{fit} and
\texttt{predict}. Does that establish that they are statistically
interchangeable? List three additional contracts that must agree.}

\section{Shared affine structure, different output semantics}

For squared-error regression, choose the identity activation and per-example
loss
\[
  a_i=z_i,
  \qquad
  \ell_i=\frac12(a_i-y_i)^2.
\]
The chain rule gives
\[
  \frac{\partial\ell_i}{\partial z_i}=a_i-y_i.
\]

For binary logistic regression, choose
\[
  a_i=\sigma(z_i)=\frac{1}{1+\exp(-z_i)}
\]
and binary cross-entropy
\[
  \ell_i=-y_i\log a_i-(1-y_i)\log(1-a_i).
\]
Since $\sigma'(z)=\sigma(z)(1-\sigma(z))$, the factors cancel:
\[
  \frac{\partial\ell_i}{\partial z_i}
  =
  \frac{\partial\ell_i}{\partial a_i}
  \frac{\partial a_i}{\partial z_i}
  =a_i-y_i.
\]
Thus both canonical models have
\[
  \nabla_w L=\frac1nX^\top(a-y),
  \qquad
  \frac{\partial L}{\partial b}
  =\frac1n\mathbf 1^\top(a-y).
\]

For a linear SVM, retain the identity postactivation, encode
$y_i\in\{-1,+1\}$, and use
\[
  J(w,b)=\frac{\lambda}{2}\lVert w\rVert_2^2
  +\frac1n\sum_i\max\{0,1-y_i z_i\}.
\]
Its score derivative is $-y_i$ inside the unit margin and $0$ beyond it, with
a subgradient choice at the kink. The equality of the regression and logistic
expressions is therefore informative but not universal. It justifies sharing
an optimization shell while keeping the objective configurable; it does not
make numerical responses, probabilities, and margin scores interchangeable.

\begin{misconceptionbox}
\textbf{Misconception:} gradient descent is identical for all single-neuron
models. \textbf{Correction:} the update mechanism is shared, but the derivative
supplied to it depends on the model and loss. The residual form above is a
special property of two canonical pairings, while hinge loss supplies a
piecewise subgradient and regularization adds $\lambda w$.
\end{misconceptionbox}

\section{Pair compatible activations and losses}

An activation and a loss jointly determine both semantics and derivatives.
One can expose them as separate low-level objects, but a safe public API should
make incompatible pairings difficult to express. A bounded teaching design
uses an output-objective contract with three operations:
\[
  \text{prediction}(z),\qquad
  \text{mean loss}(z,y),\qquad
  \text{score gradient}(z,y).
\]
The regression objective supplies identity, squared error, and $z-y$. The
probabilistic classification objective supplies sigmoid, cross-entropy, and
$\sigma(z)-y$. The margin objective supplies identity, regularized hinge loss,
and a subgradient determined by the active margin set. This third case shows
why a useful contract cannot assume every objective is smooth or every output
is a probability.

This coupling also supports stable computation. Binary cross-entropy can be
written in terms of a logit $z$ as
\[
  \ell(z,y)=\log(1+\exp z)-yz.
\]
Direct evaluation may overflow. The equivalent implementation
\[
  \ell(z,y)=\operatorname{logaddexp}(0,z)-yz
\]
remains finite over a much wider floating-point range. Numerical stability is
not an internal decoration: callers depend on it.

\begin{warningbox}
Do not apply a sigmoid and then pass its result to a loss documented to consume
logits. Conversely, do not pass logits to an interface documented to consume
probabilities. The arrays may have identical shapes while representing
different mathematical objects.
\end{warningbox}

\section{Estimator and component interfaces}

A standalone estimator owns a complete fitting procedure:
\[
\begin{array}{ll}
\texttt{fit}(X,y) & \text{learn and store fitted parameters},\\
\texttt{decision\_function}(X) & \text{return affine scores},\\
\texttt{predict}(X) & \text{return model outputs}.
\end{array}
\]
The trailing underscore convention distinguishes fitted state such as
\texttt{weights\_} from construction parameters. Prediction before fitting
should raise an actionable exception rather than return an accidental default.

A differentiable network component has a different role:
\[
\begin{array}{ll}
\texttt{forward}(u) & \text{compute and retain needed intermediates},\\
\texttt{backward}(g) & \text{propagate an upstream derivative},\\
\texttt{step}(\eta) & \text{update after gradients are assembled}.
\end{array}
\]
A hidden neuron has no observed target of its own. Calling
\texttt{fit}$(X,y)$ independently on every hidden unit would therefore train
separate models, not a neural network. The network's terminal loss supplies a
single derivative that is propagated jointly through all components.

\begin{professionalbox}
Reuse the smallest abstraction whose semantics remain valid. A network may
reuse activation functions, affine calculations, validation, initializers,
and optimizers from a single-neuron lesson. It should not reuse an estimator's
independent \texttt{fit} loop as though hidden neurons were separate prediction
problems.
\end{professionalbox}

\section{Shape, state, and failure contracts}

For a fitted single-neuron estimator, useful invariants include:
\begin{itemize}
  \item $X$ has shape $(n,p)$ with $n,p>0$;
  \item $y$ has shape $(n,)$ and aligns with row order;
  \item inputs are finite under the stated missing-value policy;
  \item targets belong to the domain required by the configured objective---for
  example, $\{0,1\}$ for logistic regression and $\{-1,+1\}$ for a linear SVM;
  \item fitted weights have shape $(p,)$;
  \item prediction uses the same feature count and ordering as fitting; and
  \item learning rate and iteration count belong to documented domains.
\end{itemize}
Type annotations make part of this structure available to editors and static
checkers. NumPy-style docstrings explain shapes, units, ordering, probability
semantics, mutation, and exceptions. Runtime checks enforce important
boundaries because Python does not enforce annotations.

Randomness is state too. Components should accept or construct a local random
generator rather than mutate NumPy's global random state. A reproducibility
record includes seed, initializer, split, preprocessing, dtype, and software
version; a seed alone is not an experiment identity.

\section{Gradient checking as a software test}

Let $J:\mathbb R^q\to\mathbb R$ be differentiable. For coordinate $j$,
\[
  \frac{\partial J(\theta)}{\partial\theta_j}
  =
  \lim_{h\to0}
  \frac{J(\theta+he_j)-J(\theta-he_j)}{2h}.
\]
For small, nonzero $h$, the right side is a central-difference approximation.
It supplies an implementation independent enough to catch a missing factor,
wrong sign, or transpose.

\begin{proposition}
If $J$ has three continuous derivatives near $\theta$, the central-difference
approximation has truncation error $O(h^2)$.
\end{proposition}

\begin{proof}
Taylor expansion in directions $e_j$ gives
\[
J(\theta\pm he_j)
=J(\theta)\pm hJ_j(\theta)
+\frac{h^2}{2}J_{jj}(\theta)
\pm\frac{h^3}{6}J_{jjj}(\xi_\pm).
\]
Subtracting cancels the constant and even-order terms. Dividing by $2h$
leaves $J_j(\theta)$ plus a term proportional to $h^2$.
\end{proof}

In floating-point arithmetic, choosing $h$ too small introduces cancellation.
Gradient tests therefore compare with both relative and absolute tolerances on
tiny deterministic examples. They do not replace mathematical derivation,
prove that the loss is scientifically meaningful, or provide an efficient
training method.

\begin{workedexample}{one shell, three models}
A team implements a \texttt{SingleNeuron} estimator that accepts an output
objective. For diabetes progression, it uses identity plus squared error. For
breast-mass classification, it uses sigmoid plus binary cross-entropy.
For a margin-based comparison, it uses identity plus regularized hinge loss
with signed labels.

The team fits scaling on each training partition only. Both models call the
same validated affine and update machinery. Tests verify prediction before fit,
feature-count mismatch, objective-specific invalid labels, extreme logits,
smooth gradients, and hinge subgradients away from the kink. Experiment reports
remain different: regression reports errors in target units, logistic
regression reports probability and decision evidence, and the SVM reports
margin scores and decisions without calling its scores probabilities. The
abstraction removes duplicate mechanism without erasing statistical meaning.
\end{workedexample}

\section{API review checklist}

Before releasing a learning component, ask:
\begin{enumerate}
  \item Are mathematical objects named and shaped explicitly?
  \item Does the loss consume logits, probabilities, or unconstrained values?
  \item Are invalid target domains rejected?
  \item Which methods require fitted state?
  \item Who owns preprocessing and threshold policy?
  \item Are randomness and numerical extremes tested?
  \item Can analytical gradients be checked on a hand-computable case?
  \item If the loss is nonsmooth, is the subgradient convention documented?
  \item Can the component be serialized without hidden notebook state?
\end{enumerate}

\section{Exercises}
\begin{enumerate}
  \item Derive the bias gradient for the squared and logistic objectives.
  Explain why the final formula agrees although the statistical models differ.
  Then derive the hinge-loss bias subgradient and identify the active set.
  \item Construct an activation--loss pairing for which
  $\partial\ell/\partial z\ne a-y$. Identify every chain-rule factor.
  \item Design a threshold adapter without changing the meaning of
  \texttt{predict}. Specify boundary behavior and tests.
  \item Add $L_2$ regularization. State whether the bias is penalized and write
  a finite-difference test that would detect the wrong policy.
  \item Explain why independently fitting hidden neurons cannot train a
  multilayer network.
  \item Explain why \texttt{predict\_proba} belongs to the logistic
  configuration but not to the uncalibrated SVM configuration.
  \item Design a pull request for one new output objective. Require a
  derivation, NumPy-style documentation, extreme-input tests, and derivative or
  subgradient evidence.
\end{enumerate}

\begin{chapterreview}
Linear regression, logistic regression, and a linear SVM share an affine
single-neuron computation. Their activations, losses, regularization, target
domains, derivatives, and outputs differ. A responsible API shares mechanism
while preserving those semantics. Standalone estimators own fitting and
prediction; differentiable components participate in a joint forward and
backward computation. Shape, fitted state, randomness, failure behavior, and
numerical stability are public contracts. Numerical checks turn derivative
formulas into independently testable claims.
\end{chapterreview}

\begin{notebooklink}
Lecture 12 notebook 00 builds and tests the single-neuron API for linear
regression, logistic regression, and a linear SVM. Its reusable implementation
lives in the course package, while notebook 01 uses the component interface to
construct a multilayer network.
\end{notebooklink}
